\documentclass[../../../include/open-logic-section]{subfiles}

\begin{document}
	
\olfileid[pt]{sth}{z}{nat}
\olsection{A escolher os números naturais}

%Let us take it for granted that the informal principle \stagesinf{}
%can be justified. It is also, though, worth commenting on the
%particular axiom we extracted from that informal principle. 

Na \olref[infinity-again]{defnomega}, definimos explicitamente
a expressão ``números naturais''. Como deve entender esta
estipulação? Não é uma afirmação metafísica, mas apenas uma decisão de
\emph{tratar} certos conjuntos como os números naturais. Já abordamos razões para pensar assim na
\olref[sfr][rel][ref]{sec}, na \olref[sfr][arith][ref]{sec} e na
\olref[sfr][infinite][dedekindsproof]{sec}. Mas podemos tornar essas
razões ainda mais nítidas.

O nosso Axioma do Infinito segue \citet{VonNeumann1925}. Mas eis outro axioma, que poderíamos ter adotado em alternativa:

\begin{defish}
\emph{Axioma do Infinito de Zermelo \citeyear{Zermelo1908Untersuchungen}.} Existe um conjunto $A$ tal que $\emptyset \in A$ e
$(\forall x \in A)\{x\} \in A$. 
\end{defish}

Se tivéssemos usado o axioma de Zermelo em vez do nosso Axioma do Infinito (inspirado em von Neumann), ter-nos-iam sido dado igualmente um conjunto infinito de Dedekind e, portanto, uma álgebra de Dedekind. Na abordagem de Zermelo, o
elemento distinguido da nossa álgebra seria de novo $\emptyset$
(o nosso sucedâneo de $0$), mas a injeção seria dada pela
aplicação $x \mapsto \{x\}$, em vez de $x \mapsto x \cup \{x\}$. A
consequência mais simples é que Zermelo trata $2$ como
$\{\{\emptyset\}\}$, ao passo que nós (com von Neumann) tratamos $2$ como
$\{\emptyset, \{\emptyset\}\}$. 

Porquê escolher um axioma do Infinito em vez do outro? A razão
prática principal é que a abordagem de von Neumann ``escala'' bastante
bem para números transfinitos. Exploraremos isto a partir do
\olref[ordinals][]{chap} em diante. Porém, na simples perspectiva de
\emph{fazer aritmética}, ambas as abordagens serviriam
igualmente bem. Por isso, se alguém lhe disser que os números naturais
\emph{são} conjuntos, a pergunta óbvia é: \emph{Que conjuntos são eles?} 

Esta pergunta precisa foi tornada célebre por \citet{Benacerraf1965}. Mas
vale a pena sublinhar que é apenas o exemplo mais famoso de um
fenómeno que já encontramos muitas vezes. O ponto essencial é este. A teoria dos conjuntos nos dá uma maneira de \emph{simular} vários
tipos de entidades ``intuitivas'': os reais, os racionais, os inteiros e os
naturais, sim; mas também pares ordenados, funções e relações.
Contudo, a teoria dos conjuntos nunca nos fornece uma escolha
\emph{única} de simulação. Há \emph{sempre} alternativas
que---de modo direto---nos teriam servido tão bem. 

\end{document}
